\documentclass[10pt,a4paper]{amsart}
\usepackage[margin=19mm]{geometry}
\usepackage{amsmath,amssymb,mathtools}
\usepackage[scr=rsfs,cal=euler]{mathalfa}
\usepackage{xfrac}
\usepackage[unicode,hidelinks,bookmarks=false]{hyperref}
\newcommand{\ie}{i.e.\ }
\newcommand{\cf}{cf.\ }
\newcommand{\resp}{respectively\ }
\newcommand{\Subsection}{\subsection}
\providecommand{\nobreakdash}{\nobreak\hbox{-}}
\setlength{\emergencystretch}{2em}
\multlinegap0pt
\usepackage{bm}
\usepackage{bbm}
\usepackage{dsfont}
\usepackage{xspace}
\usepackage{cancel}
\usepackage{mathtools}
\usepackage{amsthm}
\makeatletter
\@ifundefined{corollary}{\newtheorem{corollary}{Corollary}}{}
\@ifundefined{theorem}{\newtheorem{theorem}{Theorem}}{}
\makeatother
\def\bcS{\bar{\mathcal S}}
\def\cG{{\mathcal G}}
\def\cN{{\mathcal N}}
\def\cS{{\mathcal S}}
\def\g{{\bf g}}
\def\lp{\left (}
\def\mE{{\mathbb E}}
\def\mN{{\mathbb N}}
\def\mR{{\mathbb R}}
\def\mS{{\mathbb S}}
\def\rp{\right )}
\def\x{{\bf x}}
\title[Ground state energies of multipartite $p$-spin models -- par]{Ground state energies of multipartite $p$-spin models -- partially lifted RDT view}
\author{Mihailo Stojnic}
\date{Source: arXiv:2509.05916v1 (CC BY 4.0)}
\begin{document}
\maketitle

\noindent\emph{Source and licence.} Ground state energies of multipartite $p$-spin models -- partially lifted RDT view. Version arXiv:2509.05916v1 (CC BY 4.0), distributed under the Creative Commons Attribution 4.0 International licence, as recorded in the arXiv licence metadata for this exact version and shown on the abstract page https://arxiv.org/abs/2509.05916v1. Full rights notice: licenses/arxiv-2509.05916-CC-BY-4.0.txt.\par\smallskip

\noindent\textit{Editorial scope.} Counted unit: the Corollary whose source label is \texttt{cor:cor2} in the article (source0.tex lines 1246--1261, source label \texttt{cor:cor2}), delivered together with the article's own complete original proof of it (source0.tex lines 1263--1313, one \texttt{proof} environment of 51 lines). No mathematics is written, repaired or re-derived: statement and proof are the article's own text line for line. Required context retained uncounted: thm:thm1 (lines 953--970); eq:thm1eq2 (lines 966--969); eq:mr2 (lines 982--990); eq:mr3 (lines 1001--1014); eq:mr6 (lines 1050--1053). Every definition, display and label of those lines is retained unchanged. Omitted from this package and not redistributed: 1--952, 970--981, 991--1000, 1015--1049, 1054--1245, 1262--1262, 1314--1998. Nothing inside a retained excerpt is omitted or altered: every retained body is delivered exactly as the article's lines are written, so no omission receipt of any kind accompanies this package. Citations: the retained excerpts carry no citation command at all, so no bibliography entry of the article is transcribed and no reference is redistributed. Reference adaptation: none was needed and none was made; every reference inside the delivered unit resolves inside it, so no unresolved reference or citation is produced. Printed slips kept literal: no printed word, symbol, hypothesis or number of the counted statement or of its proof was altered. Layout and class: the article's own class is \texttt{article} with packages including inputenc, fontenc, epsf, amsmath, geometry, changepage, epsfig, amssymb, amsthm, setspace, cite, mcite, algorithmic, algorithm; the wrapper is the lane's pinned \texttt{amsart} editorial wrapper at 10pt on A4, which restates the theorem-like environments the retained text needs and adds the article's own macro definitions that the retained text needs (\texttt{\textbackslash bcS}, \texttt{\textbackslash cG}, \texttt{\textbackslash cN}, \texttt{\textbackslash cS}, \texttt{\textbackslash g}, \texttt{\textbackslash lp}, \texttt{\textbackslash mE}, \texttt{\textbackslash mN}, \texttt{\textbackslash mR}, \texttt{\textbackslash mS}, \texttt{\textbackslash rp}, \texttt{\textbackslash x}), with extra wrapper packages (bm, bbm, dsfont, xspace, cancel, mathtools). The delivered numbering is the unit's own. Assets: no retained span contains a figure environment, a graphics-inclusion command, a PDF-page inclusion, a table or a TikZ picture; no image file is copied into this package, the package declares no asset dependency and no image file is redistributed with it. Licence: version arXiv:2509.05916v1 of this article is distributed under the Creative Commons Attribution 4.0 International licence, as recorded in the arXiv licence metadata for this exact version and shown on the abstract page https://arxiv.org/abs/2509.05916v1. A full rights notice is delivered at licenses/arxiv-2509.05916-CC-BY-4.0.txt. Characters: every retained excerpt is pure ASCII, so the delivered engine, XeLaTeX with its default Latin Modern text font, reports no missing character.
\begin{theorem}
\label{thm:thm1}
  Consider  $n\in\mN$ and  $p\in 2\mN$. Let $A_{i_1,i_2,\dots,i_p}\sim \cN(0,1)$ for any choice of integers $i_j\in\{1,2,\dots,n\},1\leq j \leq p$. Also, let $g$ be a standard normal variable and for $j\in\{1,2,\dots,p\}$ let $\g^{(j)}\in \mR^{n\times 1}$ and $A^{(j)}\in \mR^{n\times n}$ be comprised of standard normals. Assume that all random variables are independent among themselves. Moreover, let $\cS^{(j)}\subseteq\mS^n, 1\leq j\leq p$ and for real scalars $c_3>0$, $c_{3,u}>0$, and $c_{3,l}>0$  set $\bar{\g}\triangleq\lp \g^{(j)}\rp_{j=1,\dots,p}$, $\bar{A}\triangleq\lp A^{(j)}\rp_{j=1,\dots,p}$, $\bcS\triangleq\lp \cS^{(j)}\rp_{j=1,\dots,p}$, and
  \begin{eqnarray}
   \label{eq:thm1eq1}
\varphi(p;\bcS,n) & = &  \max_{\x^{(j)}\in\cS^{(j)}} \sum_{i_1,i_2,\dots,i_p=1}^n A_{i_1,i_2,\dots,i_p}\prod_{j=1}^{p}\x_{i_j}^{(j)} \nonumber \\
\xi(p;\bcS,n,c_3) & = &  \frac{1}{\sqrt{n}} \frac{1}{c_3}\log \lp  \mE_A e^{c_3 \varphi(p;\bcS,n) } \rp \nonumber \\
\xi_u(p;\bcS,n,c_3) & = & \frac{1}{\sqrt{n}} \lp -\frac{c_{3}}{2}(p-1)  + \frac{1}{c_{3}} \log\lp  \mE_{\bar{\g}} e^{c_{3}\max_{\x^{(j)}\in\cS^{(j)}} \sum_{j=1}^p \lp\g^{(j)}\rp^T\x^{(j)} } \rp \rp
\nonumber \\
\xi_l(p;\bcS,n,c_3) & = &   \frac{1}{\sqrt{ n}} \frac{1}{c_{3}} \log \lp \mE_{\bar{A}} e^{ \sqrt{p}^{-1}c_{3} \max_{\x^{(j)}\in\cS^{(j)}} \sum_{j=1}^p \sum_{i_1,i_2,\dots,i_p=1}^n A^{(j)}_{i_1,i_2,\dots,i_p}\prod_{k=1}^{p}\x_{i_k}^{(j)}
 } \rp.
 \end{eqnarray}
Then
\begin{eqnarray}
   \label{eq:thm1eq2}
\xi_l(p;\bcS,n,c_{3,l})  \leq \xi(p;\bcS,n,c_{3,l}) \quad \mbox{and} \quad  \xi(p;\bcS,n,c_{3,u})   \leq  \xi_u(p;\bcS,n,c_{3,u}).
\end{eqnarray}
\end{theorem}

  \begin{eqnarray}
\label{eq:mr2}
\mE \cG (\bar{\x}^{(a_1)})\cG (\bar{\x}^{(a_2)} )   & =  & \prod_{j=1}^{p}  \sum_{i_j=1}^n   \x_{i_j}^{(j,a_1)}  \x_{i_j}^{(j,a_2)}  + p-1
=  \prod_{j=1}^{p}  \lp \x^{(j,a_1)} \rp^T\x^{(j,a_2)}  + p-1 \nonumber   \\
\mE \cG_u (\bar{\x}^{(a_1)})\cG_u (\bar{\x}^{(a_2)} )   & =  & \sum_{j=1}^p \sum_{i_j=1}^n \x_{i_j}^{(j,a_1)}\x_{i_j}^{(j,a_2)}
=
\sum_{j=1}^p
 \lp \x^{(j,a_1)} \rp^T\x^{(j,a_2)}.
  \end{eqnarray}

 \begin{eqnarray}
\label{eq:mr3}
  \mE \cG (\bar{\x}^{(a_1)})\cG (\bar{\x}^{(a_2)} )
 -
\mE \cG_u (\bar{\x}^{(a_1)})\cG_u (\bar{\x}^{(a_2)} )
 & =  &
  \prod_{j=1}^{p} a^{(j)}  + p-1
-\sum_{j=1}^p
 a^{(j)}
   \nonumber \\
& \geq &  - \prod_{j=1}^{k} a^{(j)} -\sum_{j=k+1}^p
 a^{(j)}
 + p-k  +  \prod_{j=1}^{p} a^{(j)}.
  \end{eqnarray}

\begin{eqnarray}
\label{eq:mr6}
  \mE \max_{\bar{\x}} e^{c_{3,u}\cG(\bar{\x})} &\leq &   \mE \max_{\bar{\x}} e^{c_{3,u}\cG_u(\bar{\x})},
\end{eqnarray}

\begin{corollary}
\label{cor:cor2}
Assume the setup of Theorem \ref{thm:thm1} with $\g=\g^{(j)}$, $\cS^{(j)}=\cS,1\leq j\leq p$, and
 \begin{eqnarray}
   \label{eq:cor2eq1}
 \xi_u^0(p;\cS,n,c_3) & = & \frac{\sqrt{p}}{\sqrt{n}} \lp -\frac{c_{3}}{2}(p-1)  + \frac{1}{c_{3}} \log\lp  \mE_{\g} e^{c_{3}\sqrt{p}\max_{\x\in\cS} \g^T\x } \rp \rp
\nonumber \\
\xi_l^0(p;\cS,n,c_3) & = &  \frac{\sqrt{p}}{\sqrt{n}} \frac{1}{c_{3}} \log \lp \mE_{A} e^{ c_{3}  \sqrt{n}     \zeta (p;\cS,n)  } \rp  =  \frac{\sqrt{p}}{\sqrt{n}} \frac{1}{c_{3}} \log \lp \mE_{A} e^{ c_{3} \max_{\x\in\cS }  \sum_{i_1,i_2,\dots,i_p=1}^n A_{i_1,i_2,\dots,i_p}\prod_{k=1}^{p}\x_{i_k}
 } \rp.\nonumber \\
 \end{eqnarray}
Then
\begin{eqnarray}
   \label{eq:cor2eq2}
\xi^0_l(p;\cS,n,c_{3})  \leq  \xi^0_u(p;\cS,n,c_{3}).
\end{eqnarray}
\end{corollary}

\begin{proof}
For $p\in2\mN$ the above statement follows automatically by connecting the two inequalities in (\ref{eq:thm1eq2}) with a cosmetic change $c_{3} = c_{3}\sqrt{p}$. It is useful to note that it can also be obtained separately for $p\in\mN$ by adapting the reasoning of Theorem \ref{thm:thm1}. Namely, one defines two centered Gaussian processes indexed by $\x\in\cS \subseteq \mS^n$
  \begin{eqnarray}
\label{eq:cor2mr1}
\cG^0 (\x) & = & \sum_{i_1,i_2,\dots,i_p=1}^n A_{i_1,i_2,\dots,i_p}\prod_{j=1}^{p}\x_{i_j}   + g\sqrt{p-1}  \nonumber   \\
\cG_u^0 (\x) & = &   \sqrt{p} \g^T\x .
  \end{eqnarray}
For $\x^{(a_1)}$ and $\x^{(a_2)} $ the following analogue of (\ref{eq:mr2}) can be written
  \begin{eqnarray}
\label{eq:cor2mr2}
\mE \cG^0 (\x^{(a_1)})\cG^0 (\x^{(a_2)} )   & =  & \lp  \sum_{i_j=1}^n   \x_{i_j}^{(a_1)}  \x_{i_j}^{(a_2)}  \rp^p + p-1
=  \lp \lp \x^{(a_1)} \rp^T\x^{(a_2)} \rp^p + p-1 \nonumber   \\
\mE \cG^0_u (\x^{(a_1)})\cG^0_u (\x^{(a_2)} )   & =  & p \sum_{i_j=1}^n \x_{i_j}^{(a_1)}\x_{i_j}^{(a_2)}
=
p
 \lp \x^{(a_1)} \rp^T\x^{(a_2)}.
  \end{eqnarray}
  One then notes
 \begin{eqnarray}
\label{eq:mr2a}
\mE \cG^0 (\x^{(a_1)})\cG^0 (\x^{(a_1)} )
=  \lp \lp \x^{(a_1)} \rp^T\x^{(a_1)} \rp^p  + p-1
= p = p
 \lp \x^{(a_1)} \rp^T\x^{(a_1)} =
\mE \cG^0_u (\x^{(a_1)})\cG^0_u (\x^{(a_1)} ),
  \end{eqnarray}
and, after setting $a^{0}\triangleq\lp \x^{(a_1)} \rp^T\x^{(a_2)} $, repeats line-by-line all the steps between (\ref{eq:mr3}) and  (\ref{eq:mr6}) with $a^{(j)}=a^{(0)},1\leq j\leq p$, to arrive at the following
\begin{eqnarray}
\label{eq:cor2mr6}
  \mE \max_{\x} e^{c_{3}\cG^0(\x)} &\leq &   \mE \max_{\x} e^{c_{3}\cG^0_u(\x)}.
\end{eqnarray}
A combination of (\ref{eq:cor2mr1}) and (\ref{eq:cor2mr6}) gives
\begin{eqnarray}
\label{eq:cor2mr7}
  \mE \max_{\bar{\x}} e^{c_{3}\lp \sum_{i_1,i_2,\dots,i_p=1}^n A_{i_1,i_2,\dots,i_p}\prod_{k=1}^{p}\x_{i_k}   + g\sqrt{p-1}   \rp  } &\leq &   \mE \max_{\bar{\x}} e^{c_{3} \sqrt{p} \g^T\x   },
  \end{eqnarray}
which is equivalent to
 \begin{eqnarray}
\label{eq:cor2mr9}
   \log \lp \mE e^{c_{3} \max_{\x}  \lp \sum_{i_1,i_2,\dots,i_p=1}^n A_{i_1,i_2,\dots,i_p}\prod_{k=1}^{p}\x_{i_k}    \rp  }
   \rp
   &\leq &
   -\frac{c_{3}^2}{2}(p-1) + \log \lp
    \mE e^{c_{3} \sqrt{p}  \max_{\x}  \g^T\x   }
    \rp,
\end{eqnarray}
and, after scaling by $c_3$, sufficient to complete the proof.



 \end{proof}

\end{document}
